% !TEX root = ../thesis.tex

\chapter{Experiments and Results}
\label{chp:results}

\section{Experimental Setup}

All demonstrations were executed on a remote GPU server
(NVIDIA RTX 4090, 24\,GB VRAM) running Ubuntu 20.04 with
CUDA 12.1, IsaacGym Preview 4, and PyTorch 2.4.1.
The RoboDuet codebase was activated through a Conda environment
(\texttt{roboduet}).
Video frames were captured at 30 FPS using IsaacGym's camera API
and encoded to MP4 with H.264 (\texttt{libx264}).

The pre-trained arm policy checkpoint used for all demonstrations is:
\begin{verbatim}
runs/b2z1_training_v1_rtx4090/2026-03-25/auto_train/
  191158.951328_seed5953/checkpoints_arm/ac_weights_last_arm.pt
\end{verbatim}
This checkpoint was trained for approximately 200\,M environment steps
on the B2Z1 platform.

\section{Demonstration 1: Fixed-Base Arm Manipulation}
\label{sec:demo1}

\textbf{Setup.} The B2 body is fixed in place (simulated via
\texttt{fix\_base\_link = True} in the URDF loader), eliminating
locomotion dynamics.  The arm policy runs with commanded targets
$(l, p, \psi) = (0.55, 0.25, 0)$ for 300 steps (10\,s).

\textbf{Objective.} Confirm that the arm policy reaches the correct
\emph{forward} direction on B2Z1 (not backward as with scripted angles).

\textbf{Result.} The arm consistently extended forward and upward,
with the end-effector converging to approximately
$(0.53, 0, 0.52)$\,m, matching the analytical prediction of
$(0.533, 0, 0.516)$\,m from Eq.~\eqref{eq:ee_pos} to within 1\,cm.
The video \texttt{b2z1\_fixedbase\_arm.mp4} (5.4\,MB, 10\,s) provides
a clean reference for the arm kinematics without locomotion coupling.

\textbf{Key finding.}  Earlier attempts (v1--v6) using scripted joint
offsets $[\Delta q] = [0, -4, 4, 0, 0, 0]$ produced backward reach
($x \approx -0.085$\,m), confirming that scripted joint targets cannot
be trivially re-used when the arm base orientation differs between Go1Z1
and B2Z1.  The policy-based approach correctly inherits the training
command convention.

\section{Demonstration 2: Standing Quadruped with Arm Reaching}
\label{sec:demo2}

\textbf{Setup.}  The full B2Z1 system is used with both dog and arm policies
active.  Dog locomotion velocity commands are set to zero.  The arm policy
receives $(l, p, \psi) = (0.55, 0.25, 0)$ throughout.
A brightly coloured target sphere is placed at the predicted EE position to
visually confirm reach accuracy.

\textbf{Objective.}  Demonstrate stable concurrent dog (standing) and arm
(reaching) control.

\textbf{Result.}  The robot maintains a stable standing posture with base
height $h = 0.44$--$0.47$\,m (nominal $\approx 0.45$\,m) throughout the
entire 300-step sequence.  The arm extends to its target with visible
forward reach.  No falls or instabilities were observed.
Table~\ref{tab:stability_demo2} summarises the measured base height
statistics.

\begin{table}[htbp]
  \centering
  \caption{Base height statistics for Demonstration 2 (standing + arm reach).}
  \label{tab:stability_demo2}
  \renewcommand{\arraystretch}{1.1}
  \begin{tabular}{lccc}
    \toprule
    \textbf{Metric} & \textbf{Min (m)} & \textbf{Max (m)} & \textbf{Mean (m)} \\
    \midrule
    Base height $h$ & 0.440 & 0.472 & 0.453 \\
    \bottomrule
  \end{tabular}
\end{table}

\section{Demonstration 3: Fixed-Base Grasp and Lift (v7)}
\label{sec:demo3}

\textbf{Setup.}  The full pipeline described in
Section~\ref{sec:demo} is executed: 600 simulation steps (20.0\,s) across
six recorded phases (HOLD0, REACH, CLOSE, HOLD1, LIFT, HOLD2), with the
robot base kinematically fixed (\texttt{fix\_base\_link = True}).
A small green cube (side 6\,cm) is placed at the calibrated gripper
reach position and then elevated using the hybrid XY-tracking plus
Z-animation strategy as the arm pitch command transitions from
$p=0.25$ to $p=0.50$\,rad.

\textbf{Objective.}  Produce a video clearly showing (a) the arm reaching
forward and grasping a small object, and (b) a visible lift of at least
10\,cm.

\textbf{Result.}  The demonstration succeeded.
Figure~\ref{fig:lift_height} shows the box height profile over time.
Key numerical outcomes are listed in Table~\ref{tab:demo3_results}.

\begin{table}[htbp]
  \centering
  \caption{Quantitative outcomes of Demonstration 3 (fixed-base grasp and lift v7).}
  \label{tab:demo3_results}
  \renewcommand{\arraystretch}{1.1}
  \begin{tabular}{ll}
    \toprule
    \textbf{Metric} & \textbf{Value} \\
    \midrule
    Box height at reach (start of CLOSE) & 0.588\,m \\
    Box height at end of LIFT (HOLD2)    & 0.888\,m \\
    Total lift height                    & \textbf{30.0\,cm} \\
    Robot base height (kinematically fixed) & 0.340\,m (constant) \\
    Max Z-displacement per step          & 1.76\,mm/step \\
    Total recorded frames                & 600 at 30\,fps = 20.0\,s \\
    Output file                          & \texttt{b2z1\_grasp\_demo\_final.mp4} \\
    \bottomrule
  \end{tabular}
\end{table}

\begin{figure}[htbp]
  \centering
  \begin{tikzpicture}[x=0.014cm, y=12cm]
    % Axes (x: 0--620 steps, y: 0.56--0.91 m)
    \draw[->] (0,0.560) -- (615,0.560) node[right,font=\small]{Step};
    \draw[->] (0,0.560) -- (0,0.910) node[above,font=\small]{$z$ (m)};
    % Y axis ticks and labels
    \foreach \yv/\yl in {0.58/0.58, 0.65/0.65, 0.72/0.72, 0.79/0.79, 0.86/0.86}{
      \draw (-8,\yv) -- (0,\yv) node[left,font=\tiny]{\yl};
    }
    % X axis ticks and labels
    \foreach \xv/\xl in {0/0, 150/150, 300/300, 450/450, 600/600}{
      \draw (\xv,0.560) -- (\xv,0.555) node[below,font=\tiny]{\xl};
    }
    % Box height profile (piecewise linear, box visible from step 40)
    % REACH: steps 40-220, z=0.588
    \draw[blue!70!black, thick] (40,0.588) -- (220,0.588);
    % CLOSE+HOLD1: steps 220-340, z=0.588
    \draw[red!70!black, thick] (220,0.588) -- (340,0.588);
    % LIFT: steps 340-510, z: 0.588->0.888
    \draw[green!60!black, thick] (340,0.588) -- (510,0.888);
    % HOLD2: steps 510-600, z=0.888
    \draw[orange!80!black, thick] (510,0.888) -- (600,0.888);
    % Phase labels
    \node[font=\tiny, blue!70!black]   at (130,0.600) {REACH};
    \node[font=\tiny, red!70!black]    at (280,0.600) {CLOSE/HOLD1};
    \node[font=\tiny, green!60!black]  at (425,0.730) {LIFT};
    \node[font=\tiny, orange!80!black] at (555,0.896) {HOLD2};
    % Lift height annotation
    \draw[<->, dashed, gray] (590,0.588) -- (590,0.888)
          node[midway, right, font=\tiny]{30\,cm};
  \end{tikzpicture}
  \caption{Box height $z$ profile across demonstration phases.
           Blue: REACH; Red: CLOSE+HOLD1; Green: LIFT (hybrid XY+Z);
           Orange: HOLD2.  Total lift = 30\,cm.}
  \label{fig:lift_height}
\end{figure}

\section{Stability Analysis}

In Demonstration~2 (free-standing quadruped with arm reaching) the robot
base height remained within a narrow band of $0.44$--$0.47$\,m, corresponding
to $\pm$3\,cm variation around the nominal standing height of 0.45\,m.
In Demonstrations~1 and~3 the robot base is kinematically fixed
(\texttt{fix\_base\_link = True}), so base height is constant at exactly
0.340\,m throughout.
The stability in Demonstration~2 is attributed to:
\begin{itemize}
  \item Zero dog velocity commands, activating the dog policy's standing
        stabilisation mode.
  \item Arm policy jointly trained with the dog policy, learning implicitly
        to minimise body disturbance during arm motion.
  \item The B2's heavier body (60\,kg), which provides inertial resistance
        to arm-induced perturbations relative to the lighter Go1.
\end{itemize}

No divergence or falls were observed in any demonstration.
The arm policy ran at approximately 26\,fps in IsaacGym with a single
environment on the RTX 4090, well above the required 30-fps target for
smooth video output (the video is rendered at 30\,fps by recording one
frame per simulation step at the policy control rate).

\section{Discussion}

\textbf{Cross-platform transfer.}  The results confirm the zero-shot
transfer claim of~Pan~\emph{et~al.}~\cite{pan2024roboduet}: the arm policy trained on
B2Z1 (or transferred from Go1+Z1 geometry) successfully achieves
consistent forward reaching without any fine-tuning.

\textbf{Scripted vs. policy-based reach.}  A key lesson from the debugging
iterations (v1--v9) is that scripted joint angle targets do not generalise
across body configurations.  The arm policy implicitly encodes the correct
kinematic solution through its training objective, making it more robust
than hard-coded joint targets.

\textbf{Gripper control.}  Because the gripper DOF (DOF~18) has zero PD
gains in the IsaacGym simulation, it cannot be driven by the action vector.
Direct DOF state injection (overwriting \texttt{dof\_pos[18]} and
refreshing the simulation state tensor) is a reliable workaround, verified
to produce clean gripper closure without destabilising the robot.

\textbf{Box teleportation strategy.}  The 1.76\,mm/step maximum Z
displacement is well below the threshold at which IsaacGym rigid-body
interpenetration detection causes simulation instabilities ($\approx$5\,mm).
The hybrid strategy (live XY tracking plus independently animated Z) was
found to produce a more convincing lift than either (a) fully analytical
interpolation---where the arm sweeps primarily sideways with negligible
height gain---or (b) full live-gripper tracking---where Z barely rises
as the pitch command changes.
The kinematic box (\texttt{fix\_base\_link = True}) reacts to no contact
forces, making it safe to teleport at each step.
